% subsection: 連立型

\subsection{連立型}\label{節：連立型}
  \begin{bookpoint}
    \textbf{和と差を追う}\par
    二つの式を足したり引いたりして,和と差それぞれの漸化式を作る.和と差が求まれば,連立方程式を解いて元の二つの\bookterm{sequence}{数列}を求められる.
  \end{bookpoint}
  \subsubsection{和と差で分ける}
    二つの式の係数が入れ替わっているときは,まず足し引きを試そう.
    足すと$a_n,b_n$の係数がそろい,引くと同じ大きさで符号が反対になる.
    そのため,二つの項が和や差という一つの量にまとまる.
    和と差をそれぞれ求め,元の二つの数列へ戻すところまで計算する.

    \noindent \bookblocklink{example-system-1}{problem}{answer}{【例題】}\par\nobreak
    \begin{enumerate}[label=(\arabic*),parsep=3pt,format=\bookitemlink{example-system-1}{problem}{answer}]
      \item $\displaystyle a_1=1,\,b_1=1.\ \begin{cases}a_{n+1}=3a_n+b_n\\ b_{n+1}=a_n+3b_n\end{cases}$
    \end{enumerate}

    \noindent \bookblocklink{example-system-1}{hint}{problem}{【方針】}\par\nobreak
    \begin{enumerate}[label=(\arabic*),parsep=3pt,format=\bookitemlink{example-system-1}{hint}{problem}]
      \item 係数$3,1$が入れ替わるので,和と差がそれぞれ等比型になります.和と差の初項を求めてから,$a_n,b_n$へ戻します.
    \end{enumerate}

    \noindent \bookblocklink{example-system-1}{answer}{problem}{【解答】}\par\nobreak
    \begin{enumerate}[label=(\arabic*),parsep=3pt,format=\bookitemlink{example-system-1}{answer}{problem}]
      \item $\displaystyle a_n=4^{n-1},\qquad b_n=a_n=4^{n-1}$
    \end{enumerate}

    \noindent \bookblocklink{example-system-1}{solution}{problem}{【解法】}\par\nobreak
    \begin{enumerate}[label=(\arabic*),parsep=3pt,format=\bookitemlink{example-system-1}{solution}{problem}]
      \item 二つの式を足すと$a_{n+1}+b_{n+1}=4(a_n+b_n)$なので,
            \[
              a_n+b_n=(a_1+b_1)4^{n-1}=2\cdot4^{n-1}.
                \booktag{sim-example-sum}{1}
            \]
            引くと$a_{n+1}-b_{n+1}=2(a_n-b_n)$なので,
            \[
              a_n-b_n=(a_1-b_1)2^{n-1}=0.
                \booktag{sim-example-difference}{2}
            \]
            $\bookeqbackref{sim-example-sum}{(1)}+\bookeqbackref{sim-example-difference}{(2)}$から$2a_n=2\cdot4^{n-1}$,
            $\bookeqref{sim-example-sum}{(1)}-\bookeqref{sim-example-difference}{(2)}$から$2b_n=2\cdot4^{n-1}$を得る.
            従って,$a_n=b_n=4^{n-1}$である.
    \end{enumerate}

    \bookstep{係数が対称な場合へ広げる}
    一般に,係数の位置が入れ替わった\bookterm{recurrence}{漸化式}
    \[
      \begin{cases}
        a_{n+1}=pa_n+qb_n\\
        b_{n+1}=qa_n+pb_n
      \end{cases}
    \]
    では,足すと両方の係数が$p+q$,引くと$p-q$とその符号反転になる.
    従って,和と差がそれぞれ等比型になる.
    以下の冪の式は$p+q,p-q$が非零の場合に使う.零になる方は,初項を別に保ち,$n\ge2$で$0$とする.
    \begin{align*}
      &a_{n+1}+b_{n+1}=(p+q)(a_n+b_n),\\
      &a_n+b_n=(a_1+b_1)(p+q)^{n-1}
        \booktag{sim-general-sum}{1}
    \end{align*}
    \begin{align*}
      &a_{n+1}-b_{n+1}=(p-q)(a_n-b_n),\\
      &a_n-b_n=(a_1-b_1)(p-q)^{n-1}
        \booktag{sim-general-difference}{2}
    \end{align*}
    $\bookeqbackref{sim-general-sum}{(1)}+\bookeqbackref{sim-general-difference}{(2)}$から,
    \[
      \begin{aligned}
        a_n&=\frac12\bigl\{(a_1+b_1)(p+q)^{n-1}\\
        &\qquad+(a_1-b_1)(p-q)^{n-1}\bigr\}
      \end{aligned}
    \]
    を得る.$\bookeqref{sim-general-sum}{(1)}-\bookeqref{sim-general-difference}{(2)}$からは,
    \[
      \begin{aligned}
        b_n&=\frac12\bigl\{(a_1+b_1)(p+q)^{n-1}\\
        &\qquad-(a_1-b_1)(p-q)^{n-1}\bigr\}
      \end{aligned}
    \]
    となる.次は係数が対称でない場合へ広げよう.

  \subsubsection{重みを付けて分ける}
    一般の連立\bookterm{recurrence}{漸化式}
    \[
      \begin{cases}
        a_{n+1}=pa_n+qb_n\\
        b_{n+1}=ra_n+sb_n
      \end{cases}
    \]
    でも,\bookterm{geometric}{等比数列}になる組合せを探そう.単に足すと,
    \[
      a_{n+1}+b_{n+1}=(p+r)a_n+(q+s)b_n
    \]
    となり,$p+r=q+s$でなければ和を単独で更新できない.
    そこでそれぞれに係数$\alpha,\beta$を掛けて足した,$\alpha a_n+\beta b_n$を考える.
    \[
      \begin{aligned}
        \alpha a_{n+1}+\beta b_{n+1}
          &=(p\alpha+r\beta)a_n\\
          &\quad+(q\alpha+s\beta)b_n
      \end{aligned}
    \]
    これを$\gamma(\alpha a_n+\beta b_n)$にそろえるには,係数を比較して
    \[
      \begin{cases}
        (p-\gamma)\alpha+r\beta=0\\
        q\alpha+(s-\gamma)\beta=0
      \end{cases}
    \]
    を満たす重みを選べばよい.$\alpha=\beta=0$ではすべての項が$0$となり,$a_n,b_n$を求める関係式にならないので,非零の組を探す.

    まず$\gamma$を決めよう.第1式を$s-\gamma$倍し,第2式の$r$倍を引けば$\beta$が消える.
    同様に$\alpha$を消去すると,
    \begin{align*}
      \{(p-\gamma)(s-\gamma)-qr\}\alpha&=0,\\
      \{(p-\gamma)(s-\gamma)-qr\}\beta&=0
    \end{align*}
    を得る.重みは同時に零でないので,共通の係数が零,すなわち
    \[
      \gamma^2-(p+s)\gamma+(ps-qr)=0
    \]
    が必要になる.これを連立漸化式の\bookterm{characteristic}{特性方程式}と呼ぶ.

    その解$\gamma_1,\gamma_2$を順に重みの条件へ入れ,非零の組$(\alpha,\beta),(\alpha',\beta')$を選ぶ.
    \begin{align*}
      c_n&=\alpha a_n+\beta b_n,\\
      d_n&=\alpha'a_n+\beta'b_n
    \end{align*}
    と置けば,$c_{n+1}=\gamma_1c_n,\ d_{n+1}=\gamma_2d_n$なので,倍率が非零なら,
    \[
      c_n=c_1\gamma_1^{n-1},\qquad
      d_n=d_1\gamma_2^{n-1}
    \]
    となる.初項は元の$a_1,b_1$から求める.倍率が零の場合は,その量の初項を別に保ち,$n\ge2$で$0$とする.

    元の項へは,
    \[
      \begin{cases}
        c_n=\alpha a_n+\beta b_n\\
        d_n=\alpha'a_n+\beta'b_n
      \end{cases}
    \]
    を解いて戻す.一意に復元するには,重みの二組が比例しないこと,すなわち$\alpha\beta'-\alpha'\beta\ne0$が必要である.
    \textbf{異なる二根なら,重みは比例しない.}
    もし比例すれば,同じ非零の重みが両方の係数条件を満たす.条件を引くと
    \[
      (\gamma_1-\gamma_2)\alpha=(\gamma_1-\gamma_2)\beta=0
    \]
    となり,$\gamma_1\ne\gamma_2$より$\alpha=\beta=0$.非零の重みを選んだことに矛盾する.
    異なる根は十分条件であり,必要条件ではない.
    重解では比例しない二組を取れるかで分かれる.\footnote{例えば$a_{n+1}=2a_n,\ b_{n+1}=2b_n$では$(1,0),(0,1)$を選べる.
    二組を取れない場合は,一方の量を先に求め,残る式へ代入する.\bookproblemref{practice-basic}{25}{標準問25}では,その式を指数の付加項として解く.\booknavref{節：二次特性方程式型}{隣接3項型の重解}とも比べられる.}
    このように係数を選べば,元の連立漸化式から等比型の補助数列を作れる.\footnote{非実の二根でもこの方法を使える.実数の形への復元は\sectionref{節：虚数解型隣接3項間漸化式}で扱う.}

    \noindent \hyperlink{book:example-system-1:answer:2}{\textbf{【例題】}}\par\nobreak
    \begin{enumerate}[start=2,label=(\arabic*),parsep=3pt,format=\bookitemlink{example-system-1}{problem}{answer}]
      \item $\displaystyle a_1=1,\,b_1=0,\ \begin{cases}a_{n+1}=3a_n+b_n\\ b_{n+1}=2a_n+2b_n\end{cases}$
    \end{enumerate}

    \noindent \hyperlink{book:example-system-1:problem:2}{\textbf{【方針】}}\par\nobreak
    \begin{enumerate}[start=2,label=(\arabic*),parsep=3pt,format=\bookitemlink{example-system-1}{hint}{problem}]
      \item $\gamma$を求めてから重みを選ぶと,$a_n-b_n$と$2a_n+b_n$がそれぞれ等比型になります.
    \end{enumerate}

    \noindent \hyperlink{book:example-system-1:problem:2}{\textbf{【解答】}}\par\nobreak
    \begin{enumerate}[start=2,label=(\arabic*),parsep=3pt,format=\bookitemlink{example-system-1}{answer}{problem}]
      \item $\displaystyle a_n=\frac{2\cdot4^{n-1}+1}{3},\qquad b_n=\frac{2\cdot4^{n-1}-2}{3}$
    \end{enumerate}

    \noindent \hyperlink{book:example-system-1:problem:2}{\textbf{【解法】}}\par\nobreak
    \begin{enumerate}[start=2,label=(\arabic*),parsep=3pt,format=\bookitemlink{example-system-1}{solution}{problem}]
      \item $\displaystyle a_1=1,\,b_1=0,\ \begin{cases}a_{n+1}=3a_n+b_n\\ b_{n+1}=2a_n+2b_n\end{cases}$\\
          本文の係数条件は,
          \[
            \begin{cases}
              (3-\gamma)\alpha+2\beta=0\\
              \alpha+(2-\gamma)\beta=0
            \end{cases}
          \]
          となる.非零の重みを取るための方程式は,
          \begin{align*}
            0&=(3-\gamma)(2-\gamma)-2\\
            &=\gamma^2-5\gamma+4=(\gamma-1)(\gamma-4)
          \end{align*}
          なので,$\gamma_1=1,\ \gamma_2=4$である.

          $\gamma=1$では$\alpha+\beta=0$となり,$(\alpha,\beta)=(1,-1)$を選べる.
          $c_n=a_n-b_n$と置くと,
          \[
            c_{n+1}=c_n,\quad c_1=1
            \quad\Longrightarrow\quad c_n=1
          \]
          である.$\gamma=4$では$\alpha=2\beta$となり,$(2,1)$を選べる.
          $d_n=2a_n+b_n$と置くと,
          \[
            \begin{aligned}
              &d_{n+1}=4d_n,\quad d_1=2,\\
              &\therefore\quad d_n=2\cdot4^{n-1}
            \end{aligned}
          \]
          を得る.復元は$c_n+d_n=3a_n,\ b_n=a_n-c_n$から,
          \begin{align*}
            a_n&=\frac{2\cdot4^{n-1}+1}{3},\\
            b_n&=\frac{2\cdot4^{n-1}-2}{3}
          \end{align*}
          となる.
    \end{enumerate}

    \bookstep{和と差は重みの一例}
    \bookproblemref{example-system-1}{1}{最初の例題}へ戻ろう.
    $p=s=3,\ q=r=1$を特性方程式へ入れると,
    \[
      \gamma^2-6\gamma+8=(\gamma-2)(\gamma-4)=0
    \]
    となる.$\gamma=2$なら$\alpha+\beta=0$より重み$(1,-1)$,
    $\gamma=4$なら$\alpha=\beta$より$(1,1)$を選べる.
    従って取り出す量は,
    \[
      c_n=a_n-b_n,\qquad d_n=a_n+b_n
    \]
    であり,更新は$c_{n+1}=2c_n,\ d_{n+1}=4d_n$となる.
    初項$c_1=0,\ d_1=2$も最初の解法と一致する.
    対称な係数なら足し引きが早い.重みを求める方法は,非対称な例題(2)にも使える.
    和と差も重みを付ける方法の一例なのである.
